\documentclass[10pt,a4paper]{article}
\usepackage[margin=0.8in]{geometry}
\usepackage[T1]{fontenc}
\usepackage{amsmath,amssymb,booktabs,array,tabularx,placeins,microtype,graphicx,subcaption}
\usepackage[numbers,sort&compress]{natbib}
\usepackage{xurl}
\usepackage[hidelinks]{hyperref}
\providecommand{\Description}[1]{}
\providecommand{\bibinfo}[2]{#2}
\pdfinfoomitdate=1
\pdftrailerid{}
\raggedbottom
\emergencystretch=1.5em
\hypersetup{pdfauthor={Anonymous authors},pdftitle={Jira Enhancer: Supplementary Text -- Appendices A and B},pdfsubject={Scoring checks and weighting sensitivity}}
\title{Jira Enhancer: Supplementary Text\\Appendices A and B}
\author{Anonymous authors}
\date{}
\begin{document}
\begin{center}
\Large Jira Enhancer: Supplementary Text\\Appendices A and B\\[0.4em]
\normalsize Anonymous authors
\end{center}
\appendix
\setcounter{table}{6}
\renewcommand{\thetable}{\thesection.\arabic{table}}
\section{Scoring Model Sanity Checks under Deterministic Perturbation}\label{sec:stress-testing}
Table~\ref{tab:scoring-sanity-checks} summarizes the formula-generated profiles for the internal and public-input checks. The internal check applies seeded perturbations to reused runtime-derived inputs; the public-input check applies the same formulas to 300 SWE-bench Lite issue texts~\citep{R26}. B0--B3 identify scoring profiles in this table. Their composite uses weights of 0.45 for quality, 0.30 for efficiency, and 0.25 for governance on a 0--100 scale. These rows check the scoring implementation and are separate from the executed baseline comparison in Section~5.3 of the main manuscript.

\begin{table*}[htbp]
\centering
\setlength{\tabcolsep}{3pt}
\small
\caption{Scoring model sanity checks under deterministic perturbation. All entries are formula-derived means.}
\label{tab:scoring-sanity-checks}
\begin{tabular}{@{}llrrrrr@{}}
\toprule
Input & Profile & Rows & Quality & Efficiency & Governance & Composite \\
\midrule
Seeded internal units & B0 & 5,000 & 60.17 & 30.28 & 35.18 & 44.96 \\
 & B1 & 5,000 & 59.66 & 47.01 & 36.45 & 50.06 \\
 & B2 & 5,000 & 73.76 & 52.33 & 41.97 & 59.38 \\
 & B3 & 5,000 & 88.18 & 66.97 & 68.76 & 76.96 \\
\midrule
Public issue text & B0 & 300 & 60.39 & 30.62 & 35.21 & 45.16 \\
 & B1 & 300 & 57.72 & 48.20 & 35.34 & 49.27 \\
 & B2 & 300 & 79.36 & 48.39 & 42.79 & 60.93 \\
 & B3 & 300 & 89.92 & 66.65 & 69.72 & 77.89 \\
\bottomrule
\end{tabular}
\end{table*}




\clearpage
\section{Sensitivity to Uniform Requirement-Dimension Weights}\label{sec:weight-sensitivity}
We evaluated one fixed alternative to the priority weights: the arithmetic mean of all seven requirement dimensions, with each weight set to exactly $1/7$ (approximately 14.286\%). The analysis uses the same corrected ratings from all 40 reviewers. It excludes the 160 masked-repeat ratings and averages the 1,680 primary ratings within each epic--method cell. The 14 matched epics remain the units of inference. No weight search or additional generation experiment was performed.

For each B3--baseline contrast, we enumerated all $2^{14}$ sign assignments for the two-sided paired test and used 100,000 paired epic-bootstrap resamples for percentile confidence intervals. The bootstrap seeds were 20261313 for B3--B2 and 20261314 for B3--B1, with identical resampling under both weight schemes. Holm adjustment covers the two contrasts within each scheme. As an additional multiplicity check, treating all four weighting--contrast tests as one family gives $p_{\mathrm{Holm}}=0.000488$ for every contrast.

\begin{table*}[htbp]
\centering
\footnotesize
\setlength{\tabcolsep}{3pt}
\caption{Sensitivity of the matched requirement-quality comparison to dimension weights. Scores retain the 1--5 scale. Intervals resample 14 paired epics; the reported Holm adjustment covers the two contrasts within each weighting scheme.}
\label{tab:weight-sensitivity}
\begin{tabularx}{\textwidth}{@{}ll>{\centering\arraybackslash}Xrr@{}}
\toprule
Weights & Contrast & Difference [95\% CI] & Holm $p$ & Positive epics \\
\midrule
Priority & B3--B2 & 0.542 [0.509, 0.573] & 0.000244 & 14/14 \\
Priority & B3--B1 & 1.592 [1.205, 1.984] & 0.000244 & 14/14 \\
Uniform & B3--B2 & 0.559 [0.547, 0.571] & 0.000244 & 14/14 \\
Uniform & B3--B1 & 1.547 [1.210, 1.873] & 0.000244 & 14/14 \\
\bottomrule
\end{tabularx}
\end{table*}

The uniform-weight means are 2.737, 3.724, and 4.283 for B1, B2, and B3, respectively. Table~\ref{tab:weight-sensitivity} shows positive contrasts on every epic under either specified weighting. This finding shows that the direction of comparison in the corrected dataset does not depend on the selected priority vector. It does not independently validate the rating corrections or establish stability on new portfolios.


\FloatBarrier
\begin{thebibliography}{1}
\bibitem[{Jimenez et~al.(2024)Jimenez, Yang, Wettig, Yao, Pei, Press, and
  Narasimhan}]{R26}
\bibinfo{author}{C.~E. Jimenez}, \bibinfo{author}{J.~Yang},
  \bibinfo{author}{A.~Wettig}, \bibinfo{author}{S.~Yao},
  \bibinfo{author}{K.~Pei}, \bibinfo{author}{O.~Press},
  \bibinfo{author}{K.~Narasimhan},
\newblock \bibinfo{title}{Swe-bench: Can language models resolve real-world
  github issues?},
\newblock in: \bibinfo{booktitle}{International Conference on Learning
  Representations (ICLR)}, \bibinfo{publisher}{OpenReview.net},
  \bibinfo{address}{Vienna, Austria}, \bibinfo{year}{2024}.
%Type = Article
\end{thebibliography}
\end{document}
